\documentclass[12pt]{article}
\usepackage[utf8]{inputenc}
\usepackage[T1]{fontenc}
\usepackage[margin=1in]{geometry}
\usepackage{setspace}
\usepackage{xcolor}
\usepackage[colorlinks=true, urlcolor=blue!60!black, linkcolor=blue!60!black]{hyperref}
\setlength{\parindent}{0em}
\setlength{\parskip}{0.6em}
\renewcommand{\baselinestretch}{1.0}
\begin{document}
\begin{center}
{\LARGE \textit{National Republican Senatorial Committee v. Federal Election Commission}, 609 U.S. \_\_\_ (2026)\par}
\vspace{0.5em}
{\large Case Brief --- Two-Page Edition\par}
\end{center}

\section{Memory Jogger}
Overruling \textit{Colorado II}, the Court struck down FECA's limits on parties' coordinated expenditures with their candidates as a disproportionate ``fourth prophylaxis'' (\href{https://www.supremecourt.gov/opinions/25pdf/24-621_h315.pdf}{Slip Opinion}).

\section{Detailed Case Facts}
52 U.S.C. \S~30116(d) caps party spending in ``coordination'' with federal candidates. \textit{Colorado II} (2001) upheld the caps to prevent donors from circumventing candidate contribution limits. After \textit{Citizens United}, Super PACs could spend without limit, raising over \$15.7 billion in 2023--2024 to parties' \$2.7 billion. The NRSC, the NRCC, JD Vance, and Steve Chabot sued the FEC in 2022.

\section{Procedural History}
The en banc Sixth Circuit rejected the challenge as bound by \textit{Colorado II}. The government declined to defend the law, so the Court appointed an \textit{amicus} and allowed Democratic committees to intervene. It reversed on June 30, 2026.

\section{Judicial Votes}
6--3. Kavanaugh wrote for Roberts, Thomas, Alito, Gorsuch, and Barrett; Kagan dissented, joined by Sotomayor and Jackson.

\section{Holding}
The coordinated-expenditure limits violate the First Amendment, and \textit{Colorado II} is overruled. Even under closely drawn scrutiny, a restriction must target \textit{quid pro quo} corruption and be necessary, narrowly tailored, and not disproportionate. The case was not moot.

\section{Analysis of Opinions}
\textbf{Majority:} Coordinated spending is party speech. Only \textit{quid pro quo} corruption counts, and base limits, earmarking rules, and disclosure already guard against circumvention, so the caps add an unnecessary ``fourth prophylaxis.'' \textit{Colorado II} was a ``three-legged stool where all three legs have already been knocked out'' and had weakened parties relative to Super PACs. \textbf{Kagan:} coordinated spending is ``as useful to the candidate as cash''; joint fundraising committees let one donor give \$551,300, which the party can spend on one candidate, and earmarking and disclosure cannot stop it.

\section{Examples: Future Applications}
\textbf{Same-Side:} a cap a State imposes on its own parties' coordinated spending with a gubernatorial nominee is invalid. \textbf{Opposite-Side:} the limit on Super PAC coordinated expenditures, reserved in footnote 6, likely stands. \textbf{Fence-Sitter:} the \$5,000 cap on direct party cash transfers to candidates.

\section{Critique}
Scholars such as Samuel Issacharoff welcomed the ruling as ``a welcome and overdue correction'' strengthening parties against outside groups. The Brennan Center and others argue that Congress should make such trade-offs and that the ruling lifts a donor's effective ceiling from \$7,000 to over \$550,000.

\section{Key Quotations}
\begin{quote}
``To sum up: In light of the other meaningful prophylactic measures available to the Government, and given the severe infringement on First Amendment-protected political speech that ensues from limiting a political party's spending in support of its candidates, we conclude that the political-party coordinated-expenditure limits are `disproportionate' and are not `necessary' and `narrowly tailored' for the circumvention interest it seeks to protect.''
\end{quote}
\noindent --- Justice Kavanaugh, for the Court, slip op. at 21.

\end{document}
